% !TEX root = ../../main-jos.tex
\section{结束语}

本文针对微服务分布式日志采集中高开销与低价值的矛盾，提出了一种网关驱动的分布式定向日志采集框架。主要架构贡献包括四个方面：

（1）\textbf{三层协同定向采集架构}。构建了网关预筛选、节点定向采集与中心二次过滤的三层架构，以 Sidecar 模式部署，将日志减量逻辑从后端治理前置到采集前端，形成``流量感知—定向采集—语义一致入库''的完整闭环。

（2）\textbf{网关流量驱动的动态关注清单生成}。以网关流量日志作为策略源，经过高价值过滤、URL 泛化和 Top-$K$ 筛选（算法~\ref{alg:attention}），动态识别异常与慢请求的 URL 模式，生成关注清单并经 Redis 下发到各节点。定向采集决策由实际流量驱动，无需静态配置。

（3）\textbf{定向策略三次转换算法}。通过统一的 URL 泛化函数和清单版本控制，把网关预筛选、节点定向采集和中心二次过滤贯通起来，让策略在不同层级之间保持语义一致（图~\ref{fig:transform}），避免跨层策略漂移。

（4）\textbf{资源受限下的可靠传输机制}。在 Sidecar 中引入固定缓存块和压力感知指数退避，辅以 BoltDB 本地兜底，确保内存有界传输，适应边缘资源受限的环境。

基于 Go、gRPC、Redis 和 Loki 的原型在 WSL/Docker 八节点集群上完成了端到端验证，并通过 16/32 节点规模复核与基于真实分布的多进程模拟器，补充了同频率对照、Promtail 静态过滤复现和组件消融。和同架构全量采集相比，定向模式把 Loki 入库量控制在百条量级（八节点实测 72 条 vs.\ 全量 4388$\pm$1 条，详见表~\ref{tab:compare}）。这个降幅是策略过滤与源发射频率差异的联合效果，其中策略过滤在同频率模拟下的独立降幅为 67.8\%（§6.3）。

清空统计后的规模复核显示：16 节点网关流量日志增量为 6703.33$\pm$91.78 条/轮，定向模式 Loki 入库为 48.0$\pm$5.2 条/轮；32 节点网关流量日志增量为 6692.67$\pm$121.88 条/轮，定向模式 Loki 入库为 99.0$\pm$0.0 条/轮。32 节点时关注清单达到 Top-$K{=}50$ 上限。Promtail 静态过滤复现中，本文定向策略入库 233.3 条，低于静态规则的 276.3 条。组件消融中，关闭关注清单使入库量增加 211.7\%，进一步说明动态清单是采集前端减量的核心环节。Agent CPU 绝对值均低于 0.1\%。

从 AIOps 生态的角度看，本文的工作补充了标准化框架中数据采集前端的能力。从日志故障诊断链路来看，本文降低了后端解析、异常检测和根因定位的输入规模。

未来工作按优先级排列：近期将在 Kubernetes 环境开展 64 节点云原生实测，并补充 Top-$K$ 参数敏感性分析；中期将引入 DeathStarBench、Alibaba 生产迹等真实微服务负载，验证长尾分布下的覆盖率，同时构建成本量化模型，把入库降幅转化为实际的成本节约估算；远期将探索与 OpenTelemetry 追踪上下文、eBPF 应用日志挂钩的集成，以及推进 gRPC TLS、多租户策略隔离和关注清单 JSON 规范的开源与标准化。
